\documentclass[11pt,a4paper]{ctexart}

\usepackage[margin=2.2cm]{geometry}
\usepackage{amsmath}
\usepackage{amssymb}
\usepackage{algorithm}
\usepackage{algpseudocode}
\usepackage{xcolor}
\usepackage{hyperref}

\hypersetup{
  colorlinks=true,
  linkcolor=blue,
  urlcolor=blue
}

\title{M-HIBE 空区域密钥生成算法伪代码}
\author{}
\date{}

\begin{document}
\maketitle

\section{符号说明}

设数据库为 $D=\{p_1,\ldots,p_n\}$，共有 $d$ 个维度，每个维度使用 $L$ bit 表示。
固定一个维度顺序
\[
  \pi=(\pi_1,\pi_2,\ldots,\pi_d),
\]
例如三维场景中可取 $\pi=(X,Y,Z)$。

算法的核心思想是：不枚举整个多维空间，而是沿真实数据形成的前缀父节点逐层补充空区域。
若某个父节点下的下一维范围没有数据，则生成一个空区域父密钥，并用通配符覆盖所有后续维度。

\section{通用多维离线算法}

\begin{algorithm}[H]
\caption{Offline Empty Parent Key Generation}
\begin{algorithmic}[1]
\Require Dataset $D$, dimension order $\pi$, bit length $L$, master secret key $MSK$
\Ensure Global empty parent key set $K$
\State Initialize an empty prefix index $T$
\For{each point $p \in D$}
  \For{$i = 1$ to $d$}
    \State $parent \gets$ prefixes of $p$ on dimensions $\pi_1,\ldots,\pi_{i-1}$
    \State $value \gets p_{\pi_i}$
    \State Insert $value$ into $T[parent,\pi_i]$
  \EndFor
\EndFor
\State $K \gets \emptyset$
\For{each indexed parent node $parent$ in $T$}
  \State $i \gets$ next dimension after $parent$
  \State $S \gets T[parent,\pi_i]$
  \State $E \gets \textsc{CanonicalEmptyCover}([0,2^L-1] \setminus S)$
  \For{each empty prefix $e \in E$}
    \State $pattern \gets parent \parallel e \parallel *$ for all remaining dimensions
    \State $sk \gets \textsc{M-HIBE.KeyGen}(MSK, pattern)$
    \State $K \gets K \cup \{(pattern, sk)\}$
  \EndFor
\EndFor
\State \Return $K$
\end{algorithmic}
\end{algorithm}

\section{三维展开版}

对于三维顺序 $X \rightarrow Y \rightarrow Z$，离线阶段维护两个索引：
\[
  I_X[X\text{-prefix}] = \{\text{Y values under this X-prefix}\},
\]
\[
  I_{XY}[X\text{-prefix},Y\text{-prefix}] = \{\text{Z values under this }(X,Y)\text{-prefix}\}.
\]

\begin{algorithm}[H]
\caption{3D Empty Parent Key Generation}
\begin{algorithmic}[1]
\Require 3D dataset $D$, order $X \rightarrow Y \rightarrow Z$, master secret key $MSK$
\Ensure Global empty parent key set $K$
\State Build indexes $I_X$ and $I_{XY}$ from all database points
\State $K \gets \emptyset$
\For{each $X$-prefix $p_x$}
  \State $S_Y \gets I_X[p_x]$
  \State $E_Y \gets \textsc{CanonicalEmptyCover}([0,2^L-1] \setminus S_Y)$
  \For{each empty $Y$-prefix $p_y \in E_Y$}
    \State $pattern \gets (p_x, p_y, *)$
    \State $sk \gets \textsc{M-HIBE.KeyGen}(MSK, pattern)$
    \State $K \gets K \cup \{(pattern, sk)\}$
  \EndFor
\EndFor
\For{each indexed pair $(p_x,p_y)$}
  \State $S_Z \gets I_{XY}[p_x,p_y]$
  \State $E_Z \gets \textsc{CanonicalEmptyCover}([0,2^L-1] \setminus S_Z)$
  \For{each empty $Z$-prefix $p_z \in E_Z$}
    \State $pattern \gets (p_x, p_y, p_z)$
    \State $sk \gets \textsc{M-HIBE.KeyGen}(MSK, pattern)$
    \State $K \gets K \cup \{(pattern, sk)\}$
  \EndFor
\EndFor
\State \Return $K$
\end{algorithmic}
\end{algorithm}

\section{查询阶段派生算法}

查询阶段不重新枚举数据库空间，也不重新生成全局父密钥。
服务器只选择与查询范围相交的全局空区域父密钥，并将其裁剪、派生为查询范围内的空区域密钥。

\begin{algorithm}[H]
\caption{Query-Scoped Empty Key Derivation}
\begin{algorithmic}[1]
\Require Query range $Q$, global empty parent keys $K$
\Ensure Query-scoped empty key set $K_Q$
\State $K_Q \gets \emptyset$
\For{each $(pattern, sk) \in K$}
  \State $I \gets pattern \cap Q$
  \If{$I = \emptyset$}
    \State \textbf{continue}
  \EndIf
  \State $C \gets \textsc{CanonicalCover}(I)$
  \For{each child pattern $c \in C$}
    \State $child\_sk \gets \textsc{Derive}(sk,c)$
    \State $K_Q \gets K_Q \cup \{(c, child\_sk)\}$
  \EndFor
\EndFor
\State Optionally run $\textsc{SetCover}(K_Q)$ to remove redundant keys
\State \Return $K_Q$
\end{algorithmic}
\end{algorithm}

\section{简要说明}

离线阶段的复杂度由真实数据形成的前缀结构决定，而不是由整个多维空间大小决定。
对于三维顺序 $X \rightarrow Y \rightarrow Z$，每个数据点最多贡献 $O(L)$ 个 $X$ 前缀和 $O(L^2)$ 个 $(X,Y)$ 前缀，因此索引规模约为 $O(nL^2)$。
查询阶段的成本主要取决于查询范围内空区域的可压缩覆盖数量，即最终需要派生的 query-scoped empty keys 数量。

\end{document}
